\documentclass[tikz, border=6pt]{standalone}
% Shared preamble of the standalone TikZ figures (compiled by scripts/build_figures.py).
\usepackage{amsmath, amssymb}
\usepackage{xeCJK}
\setCJKmainfont[BoldFont={Hiragino Sans W6}]{Hiragino Sans W4}
\usetikzlibrary{arrows.meta, positioning, decorations.pathreplacing, calc, patterns}
\definecolor{rod}{RGB}{70,110,170}
\definecolor{rodb}{RGB}{200,120,50}
\definecolor{okc}{RGB}{60,140,80}
\definecolor{ngc}{RGB}{190,70,60}
\definecolor{boxc}{RGB}{235,240,248}
\tikzset{
  >={Stealth[length=2.4mm]},
  every node/.style={font=\small},
  box/.style={draw, rounded corners=2pt, fill=boxc, align=flush center, inner sep=4pt},
  rodA/.style={draw, fill=rod!35},
  rodB/.style={draw, fill=rodb!35},
  note/.style={font=\footnotesize, align=flush center},
  expl/.style={font=\small, align=flush left},
  brace/.style={decorate, decoration={brace, amplitude=4pt}},
}

\begin{document}
\begin{tikzpicture}[x=1.3cm, y=1.6cm]
  \draw[->] (0,0) -- (6.3,0) node[right] {$u$};
  \draw[->] (0,0) -- (0,2.6) node[above right] {$s(u) = S(\gamma(u))$};
  \draw[thick, dashed, gray] (0.6,0) -- (0.6,2.4);
  \node[note, anchor=north] at (0.6,0) {$I_{V,N}$ の下限};
  \draw[very thick, rod, domain=0.62:6, samples=80] plot (\x, {0.85*ln(\x/0.55)});
  \pgfmathsetmacro{\uo}{2.4}
  \pgfmathsetmacro{\so}{0.85*ln(\uo/0.55)}
  \pgfmathsetmacro{\sl}{0.85/\uo}
  \draw[okc, thick] ({\uo-1.3}, {\so-1.3*\sl}) -- ({\uo+1.5}, {\so+1.5*\sl});
  \fill[okc] (\uo, \so) circle (2pt);
  \node[okc, note, anchor=south east] at ({\uo-0.1}, {\so+0.1}) {傾き $\beta = 1/T$};
  \draw[->, thick, ngc] (2.4,-0.45) -- node[below, note] {かき混ぜ (Joule)：$\gamma(u) \prec \gamma(v)$} (4.6,-0.45);
  \node[note, below] at (2.4,0) {$u$};
  \node[note, below] at (4.6,0) {$v$};
  \node[expl, anchor=west, text width=12cm] at (0,-1.25)
    {かき混ぜでエネルギーを上げられるので、エントロピー原理から $s$ は単調で、$\beta \ge 0$。\\
     傾きが緩いほど温度が高い。};
\end{tikzpicture}
\end{document}
